\documentclass[reqno]{amsart}
\usepackage[utf8]{inputenc}
\usepackage{amsmath,amssymb,amsthm}
\usepackage{hyperref}

\newtheorem{theorem}{Theorem}[section]
\newtheorem{lemma}[theorem]{Lemma}
\newtheorem{corollary}[theorem]{Corollary}
\newtheorem{proposition}[theorem]{Proposition}
\newtheorem{definition}[theorem]{Definition}
\newtheorem{remark}[theorem]{Remark}

\newcommand{\R}{\mathbb{R}}
\newcommand{\C}{\mathbb{C}}
\newcommand{\Z}{\mathbb{Z}}
\newcommand{\SO}{\mathrm{SO}}
\newcommand{\Tr}{\mathrm{Tr}}

\title[The gauge-invariant operator basis for five photons]{The multilinear gauge-invariant
operator basis for five massless photons: representation theory, invariant ring, and a complete
low-degree classification}

\author{Mart\'in Baca}
\date{}

\begin{document}

\begin{center}
\fbox{\parbox{0.85\textwidth}{\centering
\textbf{Preprint v1.0 --- not peer-reviewed.} Independently checked in-house (multiple
adversarial passes, cross-validated computations, see Section~\ref{sec:verification}), but not
yet reviewed by anyone outside this project. Feedback, corrections, and criticism are genuinely
welcome --- please open an issue at
\url{https://github.com/martinbacaia/five-photon-operator-basis} or contact the author.}}
\end{center}

\begin{abstract}
We construct the $S_5$-covariant kinematic data (Mandelstam invariants and photon polarization
data) for tree-level five-photon scattering in flat space, in generic spacetime dimension $D$
(no Gram-determinant degeneracies). We show that the space of independent Mandelstam invariants
for five massless momenta carries an \emph{irreducible} $5$-dimensional representation of $S_5$,
genuinely richer than the $n=4$ case, where $S_3$ acts by simple permutation of $(s,t,u)$. We
compute the complete invariant ring of this representation (5 primary invariants of degrees
$2$--$6$, 6 secondary invariants of degrees $0,6,7,8,9,15$) and cross-validate it exactly against
Henning, Lu, Melia \& Murayama~\cite{hlmm2017}, who solved the corresponding \emph{scalar} $n=5$
problem -- confirming our representation-theoretic machinery before extending it to spin. We then
derive the gauge-invariant polarization data structure for photons at $n=5$: each particle carries
\emph{two} independent physical parameters (not one, as at $n=4$), with closed-form expressions
for them and their inverse. We identify four families of manifestly gauge-invariant building
blocks (pairwise field-strength contractions, a linear ``singlet'', a double field-strength
``sandwich'', and -- new -- a \emph{triple} field-strength sandwich), prove a parity theorem
forcing all multilinear five-photon invariants to have odd total momentum degree, and give an
exact classification of the space of such invariants through mass dimension $11$: dimension $0$ at
degree $7$, $0$ at degrees $8$ and $10$ (parity), $3$ at degree $9$, and $16$ at degree $11$
(15 from previously known building blocks plus 1 new dimension from the triple sandwich). All
results are cross-checked by at least two independent methods, and several real errors caught and
corrected during verification are documented for transparency.
\end{abstract}

\maketitle
\tableofcontents

\section{Introduction}

Classifying the tree-level $S$-matrices of massless particles in flat space, order by order in
derivatives, is a long-standing question in the amplitudes/EFT literature. For scalars, Henning,
Lu, Melia \& Murayama (HLMM)~\cite{hlmm2017} gave a systematic method: the independent kinematic
invariants for $n$ massless scalar legs carry a representation of the permutation group $S_n$, and
the space of allowed operators at each mass dimension is the Hilbert series of the corresponding
invariant ring. Chowdhury, Gadde, Gopalka, Halder, Janagal \& Minwalla
(CGGHJM)~\cite{cgghjm2019} specialized this program to $n=4$ particles \emph{with spin} (photons
and gravitons), giving a complete, unconditional classification and parameterization of all local
four-photon and four-graviton polynomial $S$-matrices at any derivative order; they then
additionally imposed the \emph{Classical Regge Growth} (CRG) conjecture --- that any consistent
classical $S$-matrix grows no faster than $s^2$ in the Regge limit --- as a further, conjectural
physical criterion to select the subclass of these already-classified operators that is
CRG-compatible (e.g.\ isolating a four-parameter family for photons, and ruling out any $D\leq6$
polynomial correction to the Einstein graviton amplitude).

The natural next step, $n=5$, splits into two problems of very different difficulty:
\begin{itemize}
    \item The \textbf{scalar} case was solved completely by HLMM themselves~\cite{hlmm2017},
    \S 5.6.1: the relevant invariant ring, its primary/secondary invariant degrees, and Hilbert
    series are all known.
    \item The \textbf{spinning} case is left open by HLMM, who note (\S 5.4.4, in the context of
    their discussion of the operator basis at $n>4$) that rings for spinning operators are, in
    general, not Cohen--Macaulay, unlike the scalar rings (which they prove are Cohen--Macaulay for
    all $n$). This underlies the \emph{graviton} case of open question \#44 of the ``100 Open
    Questions'' survey of Strings 2024~\cite{strings2024}, which asks whether the tree-level
    Einstein $S$-matrix is the unique consistent classical $n$-graviton $S$-matrix (with no
    high-spin exchange poles), noting this is established only for $n=4$ (via CRG) and posing its
    extension to $n$-point scattering as an open exercise; the analogous $n$-photon question is not
    itself part of Q44 but is the natural companion problem, and is the one this note addresses.
\end{itemize}

This note reports our results on the spinning $n=5$ (photon) case.

\begin{remark}[Scope]
This note makes \emph{no} claim about the Classical Regge Growth conjecture at $n=5$. That
question --- what the correct generalization of the classical Regge limit to $n\geq5$ particles
even is, and whether the CRG bound holds there --- is treated separately, as an open problem with
supporting (negative) computational evidence, in a companion note~\cite{crgopennote}. Nothing in
the present note depends on Regge growth being true, false, or well-defined at $n=5$.
\end{remark}

\subsection*{Organization} Section~\ref{sec:setup} fixes conventions. Section~\ref{sec:rep} gives
the $S_5$-representation on Mandelstam invariants and its invariant ring, cross-checked against
HLMM. Section~\ref{sec:polarization} derives the photon polarization data structure at $n=5$.
Section~\ref{sec:atoms} introduces the gauge-invariant building blocks, including the new
``triple sandwich'' atom. Section~\ref{sec:parity} proves a parity theorem. Section~\ref{sec:class}
gives the classification through mass dimension 11. Section~\ref{sec:verification} summarizes the
verification methodology and documents errors caught in the process. Section~\ref{sec:discussion}
discusses scope and outlook.

\section{Setup and conventions}
\label{sec:setup}

We work with $n=5$ massless momenta $p_1,\dots,p_5$ in $D$-dimensional Minkowski space
(mostly-plus signature), $p_i^2=0$, $\sum_i p_i=0$, with $D$ generic (large enough that no
Gram-determinant relation among the $p_i$ and polarizations $\varepsilon_i$ is active --- the same
regime HLMM and CGGHJM use). Each photon carries a polarization vector $\varepsilon_i$ subject to
$\varepsilon_i\cdot p_i=0$ and the residual gauge redundancy $\varepsilon_i\to\varepsilon_i+\zeta_i
p_i$. We write $d_{ab}:=p_a\cdot p_b$ and $s_{ab}:=-2d_{ab}$.

\begin{proposition}[Standard fact]
\label{prop:fieldstrength}
Suppose $S=V^\mu(\varepsilon_i)_\mu$ is a Lorentz scalar, linear in a single $\varepsilon_i$ (with
$V$ built from any other data --- other momenta, other polarizations, and possibly other invariant
tensors such as the Levi-Civita symbol) and invariant under $\varepsilon_i\to\varepsilon_i+\zeta_ip_i$
for all $\zeta_i$. If there exists a vector $n$, not built from $\varepsilon_i$, with $p_i\cdot n\neq0$
(guaranteed here by the genericity of the kinematics fixed above), then $S$ depends on
$\varepsilon_i$ only through the field strength
$F^i_{\mu\nu}:=p_{i,\mu}\varepsilon_{i,\nu}-p_{i,\nu}\varepsilon_{i,\mu}$: explicitly,
$S=F^i_{\mu\nu}n^\mu V^\nu/(p_i\cdot n)$.
\end{proposition}

\begin{proof}
Gauge invariance forces $p_i\cdot V=0$ (from $S\to S+\zeta_i\,p_i\cdot V$). Then, using $p_i\cdot
V=0$,
\[
F^i_{\mu\nu}n^\mu V^\nu = (p_i\cdot n)(\varepsilon_i\cdot V) - (\varepsilon_i\cdot n)(p_i\cdot V)
= (p_i\cdot n)\,S,
\]
which rearranges to the claimed identity. Note that this holds regardless of what other invariant
tensors were used to build $V$; only linearity in $\varepsilon_i$ is used, which is why the claim
does \emph{not} extend to atoms with degree $\geq2$ dependence on a single $\varepsilon_i$ (e.g.\ a
symmetric form $\varepsilon_i^\mu Q_{\mu\nu}\varepsilon_i^\nu$ with $Q_{\mu\nu}p_i^\nu=0$ is gauge
invariant but generically not expressible through $F^i$ alone). Every atom used below is linear in
each polarization it depends on, so this restriction is never an issue in what follows.
\end{proof}

We build every invariant below out of Lorentz contractions of the $F^i$ and the momenta.

\section{The \texorpdfstring{$S_5$}{S5}-representation on Mandelstam invariants, and its
invariant ring}
\label{sec:rep}

\subsection{Representation}

The $\binom{5}{2}=10$ ``raw'' invariants $s_{ij}$ carry the permutation representation of $S_5$ on
the edges of the complete graph $K_5$. Momentum conservation gives, for each particle $k$, the
constraint $\sum_{j\neq k}s_{kj}=0$ --- the edge-to-vertex incidence map of $K_5$, which is
$S_5$-equivariant. The space of independent Mandelstam invariants is the \emph{kernel} of this
map: automatically an $S_5$-representation, of dimension $10-5=5$.

\begin{theorem}
\label{thm:irrep}
The $5$-dimensional representation of $S_5$ described above is irreducible.
\end{theorem}

\begin{proof}
The permutation module on (unordered) pairs of a $5$-element set decomposes, by the classical
Johnson-scheme (Gelfand pair) decomposition for $S_5$ acting on $2$-subsets, as
$M^{(3,2)}\cong V^{[5]}\oplus V^{[4,1]}\oplus V^{[3,2]}$, of dimensions $1+4+5=10$; the vertex
permutation module is $M^{(4,1)}\cong V^{[5]}\oplus V^{[4,1]}$, of dimension $5$. The
edge-to-vertex incidence map is $S_5$-equivariant and, by the dimension count above, surjective
with $5$-dimensional kernel; since its image must be an $S_5$-subrepresentation of $M^{(4,1)}$ of
the full dimension $5$, the image is all of $V^{[5]}\oplus V^{[4,1]}$, and the kernel is exactly
the complementary irreducible summand $V^{[3,2]}$. This is irreducible by construction.
\end{proof}

\begin{remark}[Independent numerical cross-check]
We additionally verified this by direct computation: the character norm
$\|\chi\|^2=\frac{1}{|S_5|}\sum_g|\chi(g)|^2$, computed from the hand-derived formula
$\chi(g)=\binom{\mathrm{fix}(g)}{2}+c_2(g)-\mathrm{fix}(g)$ (where $\mathrm{fix}(g)$ is the number
of fixed points and $c_2(g)$ the number of $2$-cycles of $g$ --- exactly the character of
$M^{(3,2)}$ minus that of $M^{(4,1)}$), equals exactly $1$ when checked against the trace of the
explicit representation matrices on all $7$ conjugacy classes of $S_5$ (values
$5,1,1,-1,1,-1,0$ on classes $1^5,21^3,2^21,31^2,32,41,5$, matching the standard character table
row for the partition $[3,2]$). A brute-force homomorphism test, $\rho(g_1)\rho(g_2)=\rho(g_1g_2)$
on $40$ random pairs, found zero failures.
\end{remark}

This is qualitatively different from $n=4$, where the analogous $2$-dimensional space $(s,t)$ is
\emph{also} irreducible (as a representation of $S_3$), but is realized by literal permutation of
$(s,t,u)$ among themselves --- no such simple description exists for $n=5$.

\subsection{The invariant ring is not free}

For $n=4$, $S_3$ acts on $(s,t)$ as a genuine reflection group, so by the
Chevalley--Shephard--Todd theorem the invariant ring is a free polynomial ring (generators in
degrees $2$ and $3$).

\begin{theorem}
\label{thm:notfree}
The ring of $S_5$-invariant polynomials on the $5$-dimensional representation of
Theorem~\ref{thm:irrep} is not a free polynomial ring.
\end{theorem}

\begin{proof}
By Chevalley--Shephard--Todd, the invariant ring is free if and only if $S_5$ acts on this
representation as a group generated by pseudo-reflections (elements with eigenvalue $1$ of
multiplicity exactly $4$). A pseudo-reflection $g\neq e$ on this $5$-dimensional representation
would have eigenvalues $(1,1,1,1,\lambda)$ with $\lambda\neq1$ a root of unity, hence
$\chi(g)=4+\lambda$; since every character value of $S_5$ is a rational integer, $4+\lambda\in\Z$
forces $\lambda=-1$ (the only non-trivial real root of unity), i.e.\ $\chi(g)=3$ is the unique
character value a pseudo-reflection could have. The character values computed in
Theorem~\ref{thm:irrep} on the $7$ conjugacy classes ($5,1,1,-1,1,-1,0$) never equal $3$ for any
non-identity class, so \emph{no} element of $S_5$ can be a pseudo-reflection in this
representation. Hence the invariant ring is not free.
\end{proof}

\begin{remark}[Independent numerical cross-check]
This was also checked directly: the exact eigenvalues (computed symbolically, no numerical
approximation) of one representative of each of the $7$ conjugacy classes of $S_5$ show that no
non-identity element has eigenvalue $1$ with multiplicity $4$ --- even transpositions, the closest
candidates, fix only a $3$-dimensional subspace ($\chi=1=a-b$ with $a+b=5$ forces $a=3$), matching
the general argument above.
\end{remark}

\subsection{Complete ring structure}

We computed the complete ring structure two independent ways.

\textbf{Molien series} (character theory) gives Hilbert-series coefficients
\[
1,\,0,\,1,\,1,\,2,\,2,\,5,\,4,\,8,\,9,\,13,\,15
\]
at degrees $0$ through $11$.

\textbf{Direct computation} via Singular's \texttt{finvar.lib} (Noether normalization together
with Kemper--Steel secondary invariants) on the same explicit $S_5$ generator matrices gives $5$
primary invariants of degrees $2,3,4,5,6$, and $6$ secondary invariants of degrees $0,6,7,8,9,15$
(of sizes $1$, $158$, $283$, $495$, $638$, $3565$ terms respectively). The reconstructed Hilbert
series
\[
H(x)=\frac{1+x^6+x^7+x^8+x^9+x^{15}}{(1-x^2)(1-x^3)(1-x^4)(1-x^5)(1-x^6)}
\]
matches the independent Molien series exactly in all $12$ degrees checked.

\begin{remark}[Validation against the literature]
\label{rem:validation}
This scalar-case ring structure is \emph{already published}: HLMM~\cite{hlmm2017}, eq.~(5.50a),
give exactly the same Hilbert series, the same primary degrees $(2,3,4,5,6)$, and the same
secondary degrees $(0,6,7,8,9,15)$, for the identical problem (five identical massless scalars,
generic $D$), using the same computational tool. We report this explicitly as a validation
checkpoint confirming our machinery before extending it to the spinning case, which is where the
genuinely new content of this note begins; the scalar ring itself is \emph{not} claimed as a new
result.
\end{remark}

\section{Photon polarization data at \texorpdfstring{$n=5$}{n=5}}
\label{sec:polarization}

\subsection{Dimensional count}

For $n=4$, the four (generic) momenta span a $3$-dimensional ``scattering plane''; each
polarization splits into a transverse part $\varepsilon_i^\perp$ (living in the
$(D-3)$-dimensional complement, acted on by the little group $\SO(D-3)$) and an in-plane part
$\varepsilon_i^\parallel$. Imposing $\varepsilon_i\cdot p_i=0$ and residual gauge freedom reduces
$\varepsilon_i^\parallel$ to a \emph{single} complex parameter $\alpha_i$ per particle
(\cite{cgghjm2019}, eq.~2.11).

Repeating this count for $n=5$: the five momenta span a $4$-dimensional scattering plane; the
transverse little group is $\SO(D-4)$, and --- after the same two reductions ---
$\varepsilon_i^\parallel$ is left with \emph{two} independent physical parameters
$(\alpha_i,\alpha_i')$ per particle.

\subsection{Explicit formula and its inverse}

Writing $w_{jk}:=p_j/s_{1j}-p_k/s_{1k}$, we derive
\[
\varepsilon_1^\parallel=\alpha_1\cdot w_{23}+\alpha_1'\cdot w_{24}+a_1\cdot p_1,
\qquad a_1 \text{ pure gauge.}
\]
Verification includes: the orthogonality identity $p_1\cdot w_{jk}=0$ for all $6$ pairs; an
explicit rank computation (exact, symbolic, treating the $p_i\cdot p_j$ as the $5$ free Mandelstam
variables) showing the ``raw'' span of the $w_{jk}$ (before quotienting by the gauge direction
$p_1$) has rank exactly $3$, with $p_1$ itself contained in that span (giving physical dimension
$3-1=2$) --- a naive first count had incorrectly predicted rank $2$, an off-by-one error caught by
comparing against the known $n=4$ case, where the analogous raw span also has rank one more than
the physical dimension. We also checked, by explicit linear solve, that the omitted pairs
$w_{25},w_{34},w_{35},w_{45}$ are each an exact linear combination of the chosen basis
$\{w_{23},w_{24},p_1\}$, confirming no information is lost or double-counted.

The inverse relation, expressing $(\alpha_1,\alpha_1')$ in terms of the field-strength contractions
$C_{jk}:=p_j\cdot F^1\cdot p_k$, is a linear $2\times2$ system whose matrix is a
\emph{non-diagonal} function of the Mandelstam invariants --- genuinely coupled, unlike the
trivial $1\times1$ case at $n=4$.

\begin{remark}[Avoiding a spurious pivot dependence]
An initial attempt to fix a single ``pivot'' particle and express all polarization data in terms
of $(\alpha_i,\alpha_i')$ explicitly ran into a real obstruction: transition functions between
different pivot choices are not always trivializable by rescaling, and the construction gives a
\emph{vector bundle of rank $2$ over the choice of pivot}, not a linear representation of $S_4$.
We do not use that construction in what follows; every result from Section~\ref{sec:atoms} onward
is built from manifestly pivot-independent, gauge-invariant contractions of the field strengths
$F^i$ directly, which sidesteps this issue entirely.
\end{remark}

\section{Gauge-invariant building blocks (``atoms'')}
\label{sec:atoms}

We build every multilinear invariant from four families of manifestly gauge-invariant
contractions of the field strengths $F^i_{\mu\nu}$ and momenta, each verified numerically with
explicit $D=6$ kinematics (gauge shifts applied to check invariance to machine precision).

\begin{center}
\small
\begin{tabular}{@{}l p{5.2cm} c c@{}}
\hline
atom & definition & pol.\ slots & mom.\ degree \\
\hline
$F^i{:}F^j$ & $F^i_{\mu\nu}F^{j,\mu\nu}$ & 2 & 2 \\[2pt]
$C^k_{ab}$ & $p_a\cdot F^k\cdot p_b$ & 1 & 3 \\[2pt]
$M^{ij}_{ab}$ & $p_a\cdot F^i\cdot F^j\cdot p_b$ (``double sandwich'') & 2 & 4 \\[2pt]
$T^{ijk}_{ab}$ & $p_a\cdot F^i\cdot F^j\cdot F^k\cdot p_b$ (``triple sandwich'', new) & 3 & 5 \\
\hline
\end{tabular}
\end{center}

$F^i{:}F^j$ admits a closed form in terms of $(\alpha,\alpha')$:
\[
F^1{:}F^2 = 2(p_1\cdot p_2)(\varepsilon_1^\perp\cdot\varepsilon_2^\perp) + \left.F^1{:}F^2\right|_\parallel,
\]
\[
\left.F^1{:}F^2\right|_\parallel = -(\alpha_1\alpha_2+\alpha_1'\alpha_2')
- Q\left[\frac{\alpha_1\alpha_2'}{2d_{13}d_{24}}+\frac{\alpha_1'\alpha_2}{2d_{14}d_{23}}\right],
\]
\[
Q = d_{12}^2+d_{12}d_{13}+d_{12}d_{14}+d_{12}d_{23}+d_{12}d_{24}+d_{13}d_{24}+d_{14}d_{23}.
\]
Notably, the coefficients of $\alpha_1\alpha_2$ and $\alpha_1'\alpha_2'$ are exactly $-1$ with no
kinematic dependence; only the cross terms carry the $Q$-dependence. The same $Q$ appears in the
inverse polarization formula of \S\ref{sec:polarization} and in the linear relations between the
six $C_{jk}$ for a fixed particle --- a consistency thread across independently-derived pieces.

$M^{ij}_{ab}$ and $T^{ijk}_{ab}$ are built by propagating a momentum through $2$ (resp. $3$) field
strengths in sequence,
\[
v \mapsto v\cdot F^m := (v\cdot p_m)\varepsilon_m - (v\cdot\varepsilon_m)p_m,
\]
representing intermediate vectors purely by their dot products with all $p_x,\varepsilon_x$
(never explicit components), which avoids any choice of basis. $T$ is new to this work; the
propagation chain was validated by checking that applying it twice reproduces the independently
verified $M$, before trusting a third application.

\section{A parity theorem}
\label{sec:parity}

\begin{theorem}
\label{thm:parity}
Every multilinear five-photon invariant (linear in each of the $5$ polarizations) built using only
the metric $\eta^{\mu\nu}$ to perform Lorentz contractions (no Levi-Civita tensor; cf.\
Remark~\ref{rem:scope}) has odd total momentum degree.
\end{theorem}

\begin{proof}
By Proposition~\ref{prop:fieldstrength}, every polarization enters through its field strength
$F^i$. Consider the complete contraction graph of the invariant: all $5$ field strengths
$F^1,\dots,F^5$, with their $10$ free indices in total, each saturated by contracting with either
another field-strength index (an ``$F$-$F$'' contraction, using $2$ indices and $0$ extra momenta)
or an external momentum (an ``$F$-momentum'' contraction, using $1$ index and $1$ extra momentum)
--- these are the only two options when contractions use only the metric. (A self-contraction
$\eta^{\mu\nu}F^i_{\mu\nu}$ of a single field strength with itself vanishes identically by
antisymmetry of $F^i$, so every nontrivial $F$-$F$ contraction pairs two distinct field strengths
$F^i,F^j$, $i\neq j$; this does not affect the count below.) If there are $x$ $F$-$F$ contractions
and $y$ $F$-momentum contractions overall, saturating all $10$ indices gives $2x+y=10$, so $y$ is
even. The total momentum degree of the invariant is $5+y$ (one intrinsic momentum per field
strength, plus $y$ external contractions), possibly times a Mandelstam multiplier of even degree
(built from $d_{ab}=p_a\cdot p_b$); since $y$ is even, the total degree is
$\equiv5\equiv1\pmod2$ regardless of the contraction pattern or multiplier used. Note this argument
is applied directly to the invariant's full contraction graph, whether or not it factors as a
product of independently-contracted pieces (each of the four atoms of Section~\ref{sec:atoms} is
one such factor, but the argument does not require this factorization).
\end{proof}

\begin{remark}[Sanity check]
The identical argument for $n=4$ gives total degree $\equiv0\pmod2$, consistent with all known
$n=4$ operators (e.g.\ $(\varepsilon_i\cdot\varepsilon_j)(\varepsilon_k\cdot\varepsilon_l)$-type)
having even degree.
\end{remark}

\begin{corollary}
\label{cor:parity}
Mass dimensions $8$ and $10$ are identically zero for $n=5$.
\end{corollary}

\section{Classification through mass dimension 11}
\label{sec:class}

\begin{remark}[Scope of the classification]
\label{rem:scope}
The classification below is complete \emph{within the family of invariants that factor as a
product of the four atoms of Section~\ref{sec:atoms}, each depending on a disjoint subset of
polarizations, contracted using only the metric} (no Levi-Civita tensor, which we have not
explored). Two distinct restrictions are implicit here, and we flag both explicitly rather than
letting the classification's scope be inferred: (i) we do not claim the four listed atoms are the
only \emph{metric-only} gauge-invariant building blocks --- e.g.\ a single connected contraction
graph touching all $5$ field strengths at once (such as a cyclic trace
$F^1{}_\mu{}^\nu F^2{}_\nu{}^\rho\cdots F^5{}_\lambda{}^\mu$) is not obviously reducible to a
product of the four atoms above, and we have not exhaustively ruled out such ``entangled''
topologies; (ii) we do not claim completeness once Levi-Civita contractions are allowed. Neither
gap affects Theorem~\ref{thm:parity} (whose proof applies directly to the full contraction graph of
any metric-only invariant, factored or not), but both are real boundaries of the degree-by-degree
dimension count in this section, stated here rather than left implicit.
\end{remark}

\begin{center}
\small
\begin{tabular}{@{}c c p{7cm}@{}}
\hline
degree & dimension & status \\
\hline
7 & 0 & proven: any candidate vanishes identically under $S_5$-symmetrization \\[2pt]
8 & 0 & Corollary~\ref{cor:parity} \\[2pt]
9 & 3 & exact; combinatorially complete (below) \\[2pt]
10 & 0 & Corollary~\ref{cor:parity} \\[2pt]
11 & 16 & exact; combinatorially complete (below) \\
\hline
\end{tabular}
\end{center}

\textbf{Degree 7.} Any degree-$7$ candidate in the known topologies vanishes identically under
$S_5$-symmetrization, by an antisymmetry mechanism: the singlet $C^k_{ab}$ is antisymmetric in
$a,b$, and some permutation in the $S_5$ orbit always swaps that pair, cancelling the sum.

\textbf{Degree 9} (dimension $3$). The only slot-partitions of the $5$ polarizations reachable at
degree $9$ with the atoms of Section~\ref{sec:atoms} are $F{:}F+F{:}F+C$ (base degree $7$, needing
a degree-$2$ Mandelstam corrector) and $F{:}F+M+C$ (degree $2+4+3=9$ exactly, no corrector) ---
$M+M+C$ is degree $11$, too high, and $T$'s minimum degree is $11$ ($T+C+C=5+3+3$). An
exact-arithmetic rank/nullspace computation (exact rational matrices, cross-checked against a
$60$-digit arbitrary-precision computation on independent kinematic samples) confirms dimension
exactly $3$, with no further direction found after exhausting all contraction variants of both
partitions.

\textbf{Degree 11} (dimension $16$). The analogous partition-counting argument gives four
topologies --- $F{:}F+F{:}F+C$ (with a degree-$4$ corrector), $F{:}F+3\times C$ (exact),
$F{:}F+M+C$ (with a degree-$2$ corrector), $M+M+C$ (exact) --- which a first pass showed span a
$15$-dimensional space. An exhaustive combinatorial sweep over all contraction variants within
these four topologies ($887$ seeds, two independent batches of $20$ random high-precision
kinematic samples each, $60$-digit precision, residuals $\sim10^{-61}$) confirmed rank exactly
$15$, ruling out an incomplete-seed-set artifact.

A fifth topology, $T+C+C$ (the new triple-sandwich atom, exact degree $5+3+3=11$), was identified
by the same slot/degree counting argument --- the \emph{only} new partition the $T$ atom opens up
at this degree --- and shown to add exactly \emph{one} further dimension, confirmed in two
independent batches of high-precision kinematic samples, and shown not to grow further after also
exhausting the other two partitions $T$ opens at degree $11$ with a corrector ($T+F{:}F$, $T+M$)
--- $1297$ total seeds tested, rank stable at $16$.

\section{Verification methodology}
\label{sec:verification}

Every numeric result above was independently re-derived by at least one alternative method before
being reported as final, and the process caught several real errors, corrected in place:

\begin{sloppypar}
\begin{itemize}
    \item A convention mismatch ($\sigma$ versus $\sigma^{-1}$, i.e.\ an accidental
    anti-homomorphism instead of a genuine group homomorphism) in an early derivation of the
    $S_5$ action, caught by a full generator-set change-of-basis test (trace/determinant matching
    alone is insufficient, since both are constant on conjugacy classes and do not distinguish an
    anti-homomorphism from a homomorphism).
    \item A silent numerical corruption from forcing exact rational coefficients to nearby
    integers, caught by comparing against exact symbolic coefficient matching instead.
    \item A homogeneity bug (a hand-typed seed of the wrong degree contaminating a rank
    computation), caught by adding a permanent degree-homogeneity guardrail to the
    rank-computation utility.
    \item A floating-point precision failure at degree eleven and above (spurious,
    sample-dependent rank estimates in double precision), diagnosed and resolved by switching to
    arbitrary-precision or exact rational arithmetic.
\end{itemize}
\end{sloppypar}

We regard this track record as informative about the reliability of the numbers reported here, but
it does not substitute for independent external verification, which we consider the necessary next
step before wider circulation of this note.

\section{Discussion and outlook}
\label{sec:discussion}

This note establishes the $S_5$-representation-theoretic and combinatorial foundation for
gauge-invariant five-photon operators through mass dimension $11$, and reports a new building
block (the triple field-strength sandwich $T$) that genuinely extends the operator basis at
degree $11$ beyond what the previously known atoms capture. All of this is independent of, and
does not bear on, whether or how the Classical Regge Growth conjecture generalizes to $n\geq5$; we
discuss that question, together with supporting (negative) computational evidence, in a companion
note~\cite{crgopennote}.

Natural continuations: (i) settle whether a fifth atom type beyond Section~\ref{sec:atoms} exists,
either via Levi-Civita contractions or via ``entangled'' metric-only topologies not factoring into
the four listed atoms (both flagged in Remark~\ref{rem:scope}), which would upgrade
Section~\ref{sec:class} from a classification within a stated family to a genuine completeness
theorem; (ii) extend the classification to degree $13$; (iii) investigate whether the ring
structure found here is Cohen--Macaulay, addressing directly the question HLMM leave open for
spinning operators at $n\geq5$.

\subsection*{Acknowledgments}
The author thanks Claude (Anthropic) for assistance with derivations, computational verification,
and manuscript preparation.

\begin{thebibliography}{9}

\bibitem{hlmm2017}
B.~Henning, X.~Lu, T.~Melia, H.~Murayama, \emph{Operator bases, $S$-matrices, and their
partition functions}, JHEP \textbf{10} (2017) 199; arXiv:1706.08520.

\bibitem{cgghjm2019}
S.~D.~Chowdhury, A.~Gadde, T.~Gopalka, I.~Halder, L.~Janagal, S.~Minwalla, \emph{Classifying and
constraining local four photon and four graviton $S$-matrices}, JHEP \textbf{02} (2020) 114;
arXiv:1910.14392.

\bibitem{strings2024}
H.~Ooguri, A.~Strominger, \emph{Strings 2024: The Future of String Theory --- 100 Open
Questions}, compiled document circulated at the Strings 2024 conference (2024).

\bibitem{crgopennote}
M.~Baca, \emph{What is the classical Regge limit for $n\geq5$ particles? An open problem, with
computational evidence from a five-photon degree-$11$ search}, in preparation.

\end{thebibliography}

\end{document}
